\documentclass{article}

\usepackage{D://Documents/Latex/Latex_Template/NotesTeXV3_test}
\usepackage{subfiles}
\usepackage{wrapfig} % Required for wrapping text around figures
\usepackage{lipsum}  % Just for generating dummy text in this example

\input{D://Documents/Latex/Latex_Template/Custom_Commands.tex}

\special{dvipdfmx:config z 0}

\begin{document}

\title{{Notes on Immersed Boundary Method}\\{\normalsize{\itshape Junior Year}}}
	\author{Thomas Li}
	\affiliation{
	Undergrauate Student at New York University Shanghai\\
	}
	\emailAdd{jl15535@nyu.edu}
	\maketitle
    \pagestyle{fancynotes}
	\newpage

\subfile{IMB_Introduction/Section_1.tex}

\newpage
\subfile{Penalty_Immersed_Boundary/section2.tex}

I will explain the principles for choosing the timestep dt without entering task planning mode as it's a conceptual question.

When choosing the timestep dt for an Immersed Boundary Method (IBM) simulation (especially one solving the Navier-Stokes equations with explicit or semi-implicit time integration), you generally need to satisfy several stability constraints. The chosen dt must be the minimum of the time steps required by each of these independent conditions.

Here are the primary principles and criteria you should refer to:

1. The Courant-Friedrichs-Lewy (CFL) Condition (Advective Stability)
This is the most fundamental restriction for any explicit fluid solver. It ensures that fluid information doesn't move more than one grid cell per timestep. For a 2D simulation, the CFL condition is: $$ \Delta t_{\text{CFL}} \le C_{\text{max}} \left( \frac{\Delta x}{|u|{\text{max}}} + \frac{\Delta y}{|v|{\text{max}}} \right)^{-1} \quad \text{or simply} \quad \Delta t_{\text{CFL}} \le C_{\text{max}} \frac{\min(\Delta x, \Delta y)}{V_{\text{max}}} $$ Where:

$V_{\text{max}}$ is the maximum expected velocity magnitude in the domain.
$C_{\text{max}}$ is the Courant number, which typically must be $\le 1.0$ for stability, but often a safer value like $0.2 \sim 0.5$ is chosen in IBM to account for explicit body forces.
2. Viscous Diffusion Stability Constraint
If your Viscous term (the diffusion term $\mu \nabla^2 \mathbf{u}$) is solved explicitly, you have a strict parabolic stability limit: $$ \Delta t_{\text{visc}} \le \frac{1}{2 \nu} \left( \frac{1}{\Delta x^2} + \frac{1}{\Delta y^2} \right)^{-1} $$ Where $\nu = \mu / \rho$ is the kinematic viscosity. (Note: If your Navier-Stokes solver uses an implicit method like Crank-Nicolson for the viscous terms, this specific restriction is bypassed, and only the CFL condition applies for the fluid).

3. Immersed Boundary Elastic Stability (The biggest restrictor)
In the Immersed Boundary Method, the fluid and the boundary are coupled explicitly. The boundary exerts elastic restoring forces (spring forces $K_{\text{mem}}$ and tether/penalty forces $K_{\text{tar}}$) on the fluid.

This explicit coupling introduces severe numerical stiffness. If $\Delta t$ is too large, an elastic restoring force will over-correct the position of a boundary point, causing it to overshoot its equilibrium position in the next step, leading to explosive oscillations (unstable blow-up).

The stability limits for these explicit Hookean springs roughly scale as: $$ \Delta t_{\text{spring}} \sim \mathcal{O}\left( \frac{\Delta x}{\sqrt{K_{\text{mem}} / \rho}} \right) $$ and for the penalty force: $$ \Delta t_{\text{penalty}} \sim \mathcal{O}\left( \sqrt{\frac{m}{K_{\text{tar}}}} \right) \quad \text{or numerically} \quad \Delta t_{\text{penalty}} \propto \frac{\Delta x^2}{K_{\text{tar}}} $$

Practical Takeaway for IBM: Because of the explicit coupling of the stiff membrane forces, the timestep is almost always heavily dominated by your boundary stiffness. When you increased $K_{\text{tar}}$ to $1000$ and $K_{\text{mem}}$ to $10$ in your latest changes, you successfully compensated by drastically dropping your dt from 0.005 to 0.0001! This is exactly the right intuition.

Recommended Workflow for tuning dt:
Start Conservative: Find a dt that works without blowing up (like your 0.0001).
Monitor maximum velocity: Keep track of your maximum $\mathbf{u}$ in the domain to ensure your CFL $\approx \frac{V_{\text{max}} \Delta t}{\Delta x}$ stays well below $0.5$.
If the simulation blows up immediately (often outputting NaN everywhere): Your $K_{\text{tar}}$ or $K_{\text{mem}}$ is too stiff for your dt. You must either lower the stiffness or decrease dt.
If the simulation blows up later in the run: This is usually a CFL violation where the fluid velocity locally accelerated too much, pushing the Courant number above $1.0$. Decrease dt.
\end{document}